\section{Routing, completion, and original-edge ownership}\label{sec:routing}
This section records the graph-theoretic interfaces used by the physical
fold. They strengthen the routing and completion framework
of~\cite{OAI2026} quantitatively and keep its original-edge interpretation.
For a graph $H$, cut expansion at least $h$ means
\[
 e_H(Z,V(H)\setminus Z)\ge h\min\{\abs Z,\abs{V(H)\setminus Z}\}
                         \quad(Z\subseteq V(H)).
\]
For a supported core, cuts are taken inside its support.

\subsection{Two edge banks and smaller-order completion}
An even core at scale $w\ge w_0$ has expansion at least
$(1023/1024)2^{sw}$. A biased edge partition, with bias $511/512$ and
deviation $1/1024$, splits it into disjoint routing and completion banks
of expansions at least
\begin{equation}\label{eq:banks}
        \frac{1044483}{1048576}\,2^{sw},\qquad
        \frac{1023}{1048576}\,2^{sw},
\end{equation}
respectively. The cut-probability bound and its logarithmic scale gate
are certified uniformly for core orders at most $2^{\rho w}$.
There is no second parity correction: the cores were already corrected
in Lemma~\ref{lem:scale}.

For completion, the sampling-and-closing construction produces at most
one cycle per core vertex and leaves a graph supported on fewer than
$\eta$ times as many vertices. Apply the smaller-order induction hypothesis
to this induced residual graph and extend the packing by the constructed
cycles. Its weight is at most $(1+C\eta)$ times the core order. The strict
support inequality, rather than a comparison of average degrees, justifies
this induction call.

A fold also leaves high-degree external centres attached to a core. The
following is the precise credit needed in this case.
\begin{lemma}[High-centre completion]\label{lem:completion}
Let a graph be supported on $W=U\cup Z$, with $U\cap Z=\varnothing$ and
$2\le\abs W$. Suppose $w\ge w_A$, $2^{sw}\le\abs U\le2^{\rho w}$,
the graph has cut expansion on $U$ at least the second quantity in
\eqref{eq:banks}, and every external vertex has even degree and at least
$d$ neighbours in $U$. Suppose also
\[
 a_w=2^{\alpha w},\quad
 \deg(v,Z)\le a_w\ (v\in U),\quad
 \sqrt{2a_w\abs U}\le d\le\sqrt{2a_w\abs U}+1.
\]
If the weighted order bound with constant $K\ge0$ holds for all graphs of
strictly smaller order than $\abs W$, then the graph has a packing with
\begin{equation}\label{eq:high-credit}
 \wt_\sigma(\mathcal P)+d/2+2a_w
            \le(1+\sigma+K\eta+1/2048)\abs U.
\end{equation}
\end{lemma}
\begin{proof}
Put $q=\abs Z$ and $v=\abs U$. Double-counting edges from $Z$ to $U$
gives $qd\le a_wv$, and $d^2\ge2a_wv$ gives $q\le d/2$.
The scale inequalities give the stronger rounding budget
\[
                    q+d/2+2a_w\le v/2048.
\]
They also allow the core expansion to extend across the high centres with
the effective expansion required for completion. The completion construction
then gives weight at most
\[
                 (v+q)+(K\eta+\sigma)v.
\]
Only vertices of $U$ need a parity charge: vertices of $Z$ already have
even degree. Add $d/2+2a_w$ and use the rounding budget to obtain
\eqref{eq:high-credit}. The proof applies on the actual support $W$;
isolated padding is removed by a cost-preserving embedding. All residual
induction calls have fewer vertices than $W$.
\end{proof}

\subsection{A single routing system for all requests}
At an active scale put
\begin{equation}\label{eq:teams}
                 Q(w)=\ceil{c(w)2^{\alpha w}}.
\end{equation}
Requests consist of an ordered pair of endpoints, a team index, a set of
forbidden interior vertices, and a nonnegative vertex weight. Endpoint
load counts occurrences, so repeated endpoints are retained.

\begin{lemma}[Grouped weighted routing]\label{lem:routing}
In an even core at scale $w\ge w_A$, suppose that each vertex has endpoint
load at most $2^{\alpha w}$, each team contains at most $Q(w)$ requests,
and each forbidden set has size at most $2Q(w)$. There are paths in its
routing bank for all requests such that distinct paths are edge-disjoint;
paths in one team have disjoint interiors; each interior avoids its
forbidden set; and for every request with weight function $\omega$,
\begin{equation}\label{eq:soft-route}
       \sum_{v\text{ in its interior}}\omega(v)
                 \le\frac\gamma{Q(w)}\sum_{v\in U}\omega(v).
\end{equation}
Their lengths satisfy the scale-dependent bound
\mcode{globalRoutingLength} used in the capacity certificates.
\end{lemma}
\begin{proof}
For the large-ratio case, write $v$ for core order, $h$ for the routing
expansion, $B=Q(w)$ for team capacity and $k_0$ for endpoint capacity.
Put $q=1/c_0$, $a=a_{\rm prune}$, $R=7519120361/500000000$ and
$L=30589490361/5000000000$. The exact rational parameters satisfy
\begin{gather*}
 1+q+aq<2a^2(q+2),\qquad 4(1+q)/a<R,\\
 r_0:=1-a-1/128-1/(R\gamma)>0,\qquad
 8/r_0<(L-1/32)^2.
\end{gather*}
The scale inequalities give $k_0\le qB$ and
$R^2L^2B^2v\le h^3$. With $\ell=\floor{L\sqrt{v/h}}$ this implies
$RB\ell\le h$.

Associate to a candidate path its edge set and the marked interior vertices
$(\text{team},v)$. These are resources in
$E(G)\sqcup(\mathcal T\times V(G))$. A path of length at most $\ell$
uses at most $2\ell$ resources. To apply the Aharoni--Haxell selection
criterion~\cite{AH2000}, it suffices to show that for every nonempty
request subfamily $I$ and every resource set $Z$ of size less than
$4\ell\abs I$, some request in $I$ has an admissible path avoiding $Z$.

Write $F$ for the edge resources of $Z$ and $Y$ for its marked-vertex
resources. Put $D=4\ell\abs I$ and $x=\abs F/D$. Adaptive pruning chooses
$z,\tau>0$ with
\[
 z+\tau=a,\qquad \frac{qx}{z}+\frac{1-x}{\tau}=\frac{1+q}{a},
 \qquad \frac{qax}{2(1+q)}<\frac{z^2}{1+z}.
\]
After deletion of $F$, pruning discards at most $\abs F/(zh)$ vertices
and leaves expansion $(1-z)h$. Endpoint and team loads imply
\[
          \frac{k_0\abs F}{zh}+\frac{B\abs Y}{\tau h}<\abs I.
\]
Thus a request remains whose endpoints survive and whose team has at most
$\tau h$ marked interior obstructions. Delete those marks together with
its hard ports, and construct a short path in the surviving expanding core.
The short-path fan and averaging argument gives the additional weight bound
$(R\gamma\ell/h)\sum\omega\le(\gamma/B)\sum\omega$.
The retained expansion $r_0h$ gives the required path length $\ell$.
This proves the resource-avoidance condition. Aharoni--Haxell then selects
all paths simultaneously with disjoint resources, giving both forms of
disjointness in the statement. The formal selection lemma keeps repeated
endpoint occurrences as separate requests.

The actual integer $Q(w)$ is used throughout. Its rounding inequalities
prove the endpoint, team, hard-port and cubic-capacity gates at $w_A$ and
extend them to the full scale half-line.

There are two size regimes. If core order divided by routing-bank expansion
is at least $4096$, use the grouped weighted router directly. Otherwise
use the comparison bank and its bounded-ratio routing theorem. The
comparison ratio remains at least $4$ after the original parity correction.
Both regimes have the same conclusion~\eqref{eq:soft-route}. Thus no
intermediate range of core orders is omitted.

The completion machinery additionally uses an indexed reservoir router.
Its explicit gates are $z\ge8192$, $8192\ln v\le z$,
$2112\abs I\le z^2$, minimum degree at least $15z/16$, expansion at least
$z/30$, and endpoint load at most $z/64$. A balanced bank split leaves
one expanding bank to connect endpoint occurrences to a weak clique
immersion in the other. The clique-immersion theorem of~\cite{GLW2019}
and the terminal-pairing construction of~\cite{GMM2018}
then produce actual edge-disjoint paths. Only edge-disjointness is needed
from this reservoir theorem; it does not assert disjoint interiors.
All these graph constructions, including their integer capacities, are
proved in the formal source, not supplied by numerical gate checks alone.
\end{proof}

The key bookkeeping object for applying Lemma~\ref{lem:routing} is an
\emph{original endpoint certificate}. It assigns to each occurrence of
an endpoint of a request an incidence $(e,v)$ of an original edge, where
$v$ is that endpoint. The assignment is injective. If the original edges
used by two families are disjoint, their certificates combine. Reordering
the incoming endpoints within the same owner core preserves the certificate.

Three families are combined before routing: low two-edge paths, requests
from paired matching edges, and the switches of quotient cycles. Their
original edge sets are disjoint. The combined injection implies that the
\emph{total} endpoint load at $v$ is bounded by its degree in the assigned
descendant graph. Active-block preparation bounds this degree by
$2^{\alpha w}$. Bounding the three families separately and adding those
bounds would not establish the required load.

The combined team construction also preserves the $Q(w)$ capacity and
$2Q(w)$ hard-port limit. Apply Lemma~\ref{lem:routing} once to each selected
core, and use those same paths in all subsequent closures. Low groups and
paired matching groups close into actual simple cycles. Quotient cycles
are lifted through their switches, with the soft charge
\eqref{eq:soft-route}. The untouched completion bank remains available;
the core remainder deletes exactly the paths actually used. These ownership
rules are the reason that the four families in the next section have both
an exact edge cover and pairwise disjointness.
